%% ========================================================================
%% Brāhmasphuṭasiddhānta Part 3 — Chunk 3 (Pages 1957–2096)
%% Author: Brahmagupta (c. 598–668 CE)
%% Commentary: Hindi translation and explanation
%% Pages: 281–420 of the PDF (original pages २२–१६०)
%% ========================================================================

\documentclass[11pt,a4paper]{article}

%% -------- Core Packages --------
\usepackage{fontspec}
\usepackage{polyglossia}
\usepackage{amsmath,amssymb,amsthm}
\usepackage{geometry}
\usepackage{graphicx}
\usepackage{xcolor}
\usepackage{titlesec}
\usepackage{fancyhdr}
\usepackage{enumitem}
\usepackage{array,booktabs,longtable,multirow}
\usepackage{hyperref}
%\usepackage{verse}  % not available, using custom env
%\usepackage{lettrine} % not available
%\usepackage{microtype} % not available
%\usepackage{tocloft} % not available
\usepackage{indentfirst}
\usepackage{tikz}
%\usepackage{framed} % not available

%% -------- Page Geometry --------
\geometry{
  paperwidth=210mm,paperheight=297mm,
  left=22mm,right=22mm,top=28mm,bottom=28mm,
  headheight=14pt,footskip=30pt
}

%% -------- Language Setup --------
\setmainlanguage{english}
\setotherlanguage{sanskrit}
\setotherlanguage{hindi}

%% -------- Font Configuration --------
\setmainfont{TeX Gyre Termes}[Renderer=Harfbuzz]
\setsansfont{TeX Gyre Heros}[Renderer=Harfbuzz]
\setmonofont{DejaVu Sans Mono}[Scale=MatchLowercase]

%% Devanagari (Sanskrit/Hindi) Font
\newfontfamily\devanagarifont[Script=Devanagari,Scale=1.15,
  BoldFont={* Bold}]{Noto Serif Devanagari}
\newfontfamily\devanagarifontsf[Script=Devanagari,Scale=1.1]{Noto Sans Devanagari}

%% -------- Colors --------
\definecolor{titlecolor}{RGB}{45,55,72}
\definecolor{sectioncolor}{RGB}{55,65,81}
\definecolor{notecolor}{RGB}{120,53,15}
\definecolor{rulecolor}{RGB}{180,160,140}
\definecolor{versecolor}{RGB}{30,40,60}
\definecolor{commentarycolor}{RGB}{20,30,50}
\definecolor{shlokabg}{RGB}{250,248,245}

%% -------- Section Formatting --------
\titleformat{\section}
  {\normalfont\Large\bfseries\color{sectioncolor}}
  {\thesection}{1em}{}
\titleformat{\subsection}
  {\normalfont\large\bfseries\color{sectioncolor}}
  {\thesubsection}{1em}{}
\titleformat{\subsubsection}
  {\normalfont\normalsize\bfseries\color{sectioncolor}}
  {\thesubsubsection}{1em}{}

%% -------- Header/Footer --------
\pagestyle{fancy}
\fancyhf{}
\fancyhead[L]{\small\textit{Brāhmasphuṭasiddhānta}}
\fancyhead[R]{\small\thepage}
\renewcommand{\headrulewidth}{0.4pt}
\renewcommand{\footrulewidth}{0pt}

%% -------- Custom Commands --------
\newcommand{\longline}{\noindent\rule{\textwidth}{0.4pt}}
\newcommand{\shortline}{\noindent\rule{0.5\textwidth}{0.4pt}\par}
\newcommand{\cnote}[1]{\textcolor{notecolor}{\textit{#1}}}

%% Sanskrit/Hindi inline helpers - using long def for multiline support
\long\def\dev#1{{\devanagarifont #1}}
\newcommand{\devsf}[1]{{\devanagarifontsf #1}}

%% Verse environment for shlokas
\newenvironment{shloka}
  {\begin{center}\begin{minipage}{0.9\textwidth}\begin{center}\devanagarifont\color{versecolor}}
  {\end{center}\end{minipage}\end{center}}

%% Commentary environment
\newenvironment{commentary}
  {\begin{quote}\devanagarifont\small\color{commentarycolor}}
  {\end{quote}}

%% Section heading in Devanagari
\newcommand{\devsection}[1]{\section*{\dev{#1}}\addcontentsline{toc}{section}{#1}}
\newcommand{\devsubsection}[1]{\subsection*{\dev{#1}}}

%% Fancy chapter title box
\newcommand{\chapterbox}[1]{%
\begin{center}
\begin{tikzpicture}
  \node[draw=rulecolor, line width=1pt, inner sep=10pt, rounded corners=3pt, fill=shlokabg] {%
    \Large\bfseries\devanagarifont\color{titlecolor} #1
  };
\end{tikzpicture}
\end{center}
\vspace{1em}
}

%% -------- Hyperref Setup --------
\hypersetup{
  colorlinks=true,
  linkcolor=sectioncolor,
  citecolor=sectioncolor,
  urlcolor=sectioncolor,
  pdfencoding=auto,
  pdftitle={Brāhmasphuṭasiddhānta Part 3 — Chunk 3},
  pdfauthor={Brahmagupta}
}

%% -------- TOC Depth --------
\setcounter{tocdepth}{3}
\setcounter{secnumdepth}{3}

%% ========================================================================

\begin{document}

%% ========================================================================
%% TITLE PAGE
%% ========================================================================
\begin{titlepage}
\centering
\vspace*{2cm}

{\Huge\bfseries\color{titlecolor}\dev{ब्राह्मस्फुटसिद्धान्त}}\\[0.8cm]
{\LARGE\bfseries\color{sectioncolor}\dev{भाग तृतीय — खण्ड ३}}\\[1.5cm]

{\large\textit{Brāhmasphuṭasiddhānta}}\\[0.3cm]
{\large Part 3 — Chunk 3}\\[2cm]

\shortline

{\Large\textbf{Brahmagupta}}\\[0.5cm]
{\large c. 598–668 CE}\\[2cm]

\shortline

{\large Original pages: \dev{२२ — १६०}}\\[0.3cm]
{\large PDF pages: 281 — 420}\\[2cm]

\vfill

{\small Typeset from scanned manuscript with Sanskrit verses and Hindi commentary}\\[0.3cm]
{\small Covers chapters १७–२१ of the Grahagaṇita and Golādhyāya sections}

\end{titlepage}

%% ========================================================================
%% TABLE OF CONTENTS
%% ========================================================================
% \tableofcontents  % requires tocloft package
\newpage

%% ========================================================================
%% INTRODUCTION
%% ========================================================================
\section*{Introduction}
\addcontentsline{toc}{section}{Introduction}

This volume contains a major portion of Brahmagupta's \textit{Brāhmasphuṭasiddhānta}, one of the most important works in the history of Indian mathematics and astronomy. Written in 628 CE, this treatise consists of 24 chapters covering:

\begin{itemize}[leftmargin=*]
  \item \textbf{Grahagaṇitādhyāya} (Chapters 1–17): Planetary computations, including mean and true longitudes, conjunctions, eclipses, and other astronomical phenomena.
  \item \textbf{Kuṭṭakādhyāya} (Chapter 18): Algebraic methods for solving linear Diophantine equations (the ``pulverizer'' method).
  \item \textbf{Chāyādhyāya} (Chapter 20): Shadow computations using gnomons.
  \item \textbf{Golādhyāya} (Chapter 21): Spherical astronomy and cosmology.
\end{itemize}

The text presented here includes:
\begin{enumerate}
  \item Sanskrit verses (\textit{ślokas}) from the original text
  \item Hindi commentary (\textit{ṭīkā}) with detailed explanations
  \item Mathematical working with step-by-step derivations
\end{enumerate}

\newpage

%% ========================================================================
%% VISHAYANUKRAMANIKA — TABLE OF CONTENTS
%% ========================================================================
\section*{\dev{विषयानुक्रमणिका}}
\addcontentsline{toc}{section}{विषयानुक्रमणिका (Table of Contents)}

\chapterbox{\dev{विषयानुक्रमणिका}}

\subsection*{\dev{१७. सृङ्गोत्स्त्युतराध्यायः}}
\dev{%
\begin{tabular}{@{}p{0.8\textwidth}r@{}}
प्रश्नकथनम् & ११३६ \\
परिलेखकथनम् & ११३६ \\
प्रकारान्तरेण परिलेखकथनम् & ११४२ \\
फलके परिलेखकथनम् & ११४४ \\
विशेषकथनम् & ११४५ \\
अध्यायोपसंहारः & ११४६ \\
\end{tabular}
}

\subsection*{\dev{१८. कुट्टकाध्यायः}}
\dev{%
\begin{tabular}{@{}p{0.8\textwidth}r@{}}
कुट्टकार्थप्रयोजनम् & ११४६ \\
कुट्टकदोषाणां प्रशंसाकथनम् & ११४६ \\
कुट्टककथनम् & ११५० \\
कुट्टके विशेषकथनम् & ११५५ \\
भगणादिविशेषतो दर्शनानयनम् & ११५६ \\
विशेषकथनम् & ११५८ \\
स्थिरकुट्टककथनम् & ११६१ \\
स्थिरकुट्टकदर्शनागतिकथनम् & ११६३ \\
स्थिरकुट्टके विशेषकथनम् & ११६६ \\
विलोमगतिकथनम् & ११७१ \\
प्रश्नकथनम् & ११७२ \\
अन्यप्रश्नकथनम् & ११७४ \\
अन्य प्रश्नकथनम् & ११७५ \\
अन्य प्रश्नकथनम् & ११७६ \\
पूर्वप्रश्नस्योत्तरकथनम् & ११७७ \\
अपरप्रश्नकथनम् & ११७७ \\
पूर्वप्रश्नस्योत्तरकथनम् & ११७७ \\
\end{tabular}
}

\dev{%
\begin{tabular}{@{}p{0.8\textwidth}r@{}}
विशेषकथनम् & ११७९ \\
अन्यप्रश्नकथनम् & ११८० \\
पूर्वप्रश्नस्योत्तरकथनम् & ११८१ \\
अन्यप्रश्नकथनम् & ११८२ \\
अन्यप्रश्नकथनम् & ११८४ \\
अन्यप्रश्नकथनम् & ११८५ \\
अन्यप्रश्नकथनम् & ११८७ \\
अन्यप्रश्नकथनम् & ११८९ \\
धनर्णशून्यानां संकलनम् & ११९१ \\
व्यवकलनम् & ११९६ \\
गुणने करणसूत्रकथनम् & ११९७ \\
भागहारे करणसूत्रकथनम् & ११९९ \\
संक्रमणविभमकर्मकथनम् & ११९८ \\
समद्विबाहुत्रिभुजे भुजयोर्बर्गेणम् & ११९९ \\
करणीयोगान्तरे गुणनकथनम् & १२०१ \\
करणभागहारे करणसूत्रकथनम् & १२०१ \\
करणमूलानयनार्थं कथनम् & १२०२ \\
अव्यक्तसंकलित यावद् इष्टतयोः करणसूत्रकथनम् & १२०४ \\
अव्यक्तगुणने सूत्रकथनम् & १२०५ \\
अव्यक्तमानानयनार्थं कथनम् & १२०७ \\
वर्गसमीकरणाकथनम् & १२०७ \\
वर्गसमीकरणोद्धरणमानयनम् & १२०८ \\
प्रश्नकथनम् & १२११ \\
अन्यप्रश्नकथनम् & १२१२ \\
अन्यप्रश्नकथनम् & १२१२ \\
अन्यप्रश्नकथनम् & १२१४ \\
अन्यप्रश्नकथनम् & १२१४ \\
अनेकवर्गसमीकरणम् & १२१७ \\
प्रश्नकथनम् & १२१९ \\
अन्यप्रश्नकथनम् & १२२१ \\
अन्यप्रश्ननिरूपणम् & १२२२ \\
अन्यप्रश्नद्वयवर्णनम् & १२२४ \\
अन्यप्रश्नकथनम् & १२२५ \\
अन्यप्रश्नद्वयस्थापनम् & १२२७ \\
अन्यप्रश्नवर्णनम् & १२२८ \\
अन्यप्रश्नकथनम् & १२३१ \\
\end{tabular}
}

\newpage

\subsection*{\dev{१९. शङ्कुच्छायादिज्ञानाध्यायः}}
\dev{%
\begin{tabular}{@{}p{0.8\textwidth}r@{}}
प्रथमप्रश्नकथनम् & १२६५ \\
अन्यप्रश्नस्थापनम् & १२६५ \\
अन्यप्रश्ननिरूपणम् & १२६६ \\
अन्यप्रश्नस्थापनम् & १२६७ \\
अन्यप्रश्नकथनम् & १२६८ \\
प्रश्नान्तरस्थापनम् & १२६९ \\
प्रथमप्रश्नस्योत्तरकथनम् & १२७७ \\
उदयान्तरमस्तान्तरघटिकाद्यरितिप्रश्नद्वयस्योत्तरकथनम् & १३०१ \\
मध्यगते रन्तरसाधनं तदुत्तरकथनं च & १३०२ \\
रविवशाङ्कुमाननिरूपणम् & १३०२ \\
दीपशिखोच्चयच्छुक्लान्तरभूमिज्ञाने छायानयनस्योत्तरकथनम् & १३०४ \\
छाया द्वितीयमात्रान्तरविधानेन त्रिप्रश्नोत्तरकथनम् & १३०५ \\
छायातो गृहादीनां च्छायानयनम् & १३०७ \\
इष्टगुहोच्चयको यदिति प्रश्नस्योत्तरसम्पादनम् & १३०८ \\
गृहपुरुषान्तरसलिले यो हटचैत्यादि प्रश्नोत्तरकथनम् & १३१० \\
बीजय गुहायां सलिले प्रसारे ति प्रश्नोत्तरकथनम् & १३११ \\
उच्छ्रितिकथनम् & १३१४ \\
\end{tabular}
}

\subsection*{\dev{२०. छद्दिच्छायाध्यायः}}
\dev{%
\begin{tabular}{@{}p{0.8\textwidth}r@{}}
& १३१६--१३२० \\
\end{tabular}
}

\subsection*{\dev{२१. गोलाध्यायः}}
\dev{%
\begin{tabular}{@{}p{0.8\textwidth}r@{}}
आरम्भप्रयोजनकथनम् & १३२३ \\
भूगोलसंस्थानकथनम् & १३२३ \\
दैवासुरसंस्थानवर्णनम् & १३२४ \\
देवदैत्ययोर्मेषचक्रमरगत्यवस्थापनम् & १३२७ \\
चक्रभ्रमणव्यवस्थाकथनम् & १३२९ \\
दैनानीनां रविभ्रमणार्थदिवसानाम् & १३२९ \\
देवदैत्ययोः राशिसंस्थानम् & १३३० \\
देवदैत्ययोः पितुमानस्यैव दिनप्रमाणानिरूपणम् & १३३० \\
भूगोल लग्नावलयोः संस्थानम् & १३३४ \\
निरक्षदेशान्तरयोजनकथनम् & १३३६ \\
क्षुद्रद्योः कथानिरूपणम् & १३३७ \\
\end{tabular}
}

\newpage

%% ========================================================================
%% CHAPTER 17: GRAHAGANITADHYAYA
%% ========================================================================
\devsection{सृङ्गोत्स्त्युतराध्यायः}

\begin{center}
\textit{(Original pages ११३६–११४६; PDF pages 281–292)}
\end{center}

\longline

\devsubsection{प्रश्नकथनम्}

\begin{shloka}
द्विजार्कसनवरसेन्दूर्जिनतिथिविषयाग्रहार्थचापानाम् ।\\
अर्धज्याखण्डानि ज्यामुकैर्नैक्यं स भोग्यफलम् ॥१५॥\\[0.5em]
गतभोग्यखण्डकान्तर दलविकलविधा च्छुतैर्वर्भिराप्तैः ।\\
तच्छृतिदलं युतेन भोग्यादुनार्धिकं भोग्यम् ॥१६॥
\end{shloka}

\begin{commentary}
\dev{%
यह आचार्योंक्त भोग्यखण्ड स्पष्टीकरण है। भास्कराचार्य ने इसी को सिद्धान्तशिरो-\break
मणि के स्पष्टीकार में ``यातीतययोः खण्डकयोर्विशेषः शेषशनिर्ध्नः'' इत्यादि द्वारा लिखा\break
है। भास्कराचार्य ने १२० त्रिज्या ग्रहण की है। यहीं आचार्य ब्रह्मगुप्त ने १५० त्रिज्या\break
ग्रहण की है। यहीं यह बात बड़ी विचित्र लगती है कि आचार्योंक्त विषयों को ही सिद्धान्त\break
शेखर में श्रीपति ने सर्वत्र श्लोकान्तरों में लिखा है, परन्तु पता नहीं क्यों उन्होंने आचार्योंक्त\break
इस प्रसिद्ध भोग्य खण्ड स्पष्टीकरण की चर्चा तक नहीं की।
}
\end{commentary}

\devsubsection{पञ्चसिद्धान्तिका उद्धरण}

\begin{shloka}
रविवसतिगतः पञ्चयुग्मं वपांशि पितामहोद्दिष्टानि ।\\
अर्धमासार्द्धयातद्विमर्मसैरमर्मस्नपष्टयाऽऽत्म ॥३॥
\end{shloka}

\begin{commentary}
\dev{%
ब्राह्मस्फुट सिद्धान्त में जिन आचार्यों के नाम आते हैं। उनके सम्बन्ध में कुछ विचार\break
करते हैं। सब सिद्धान्तों में प्राचीन या सबसे प्राचीन सिद्धान्त ब्रह्मसिद्धान्त ही है। इसी को\break
लोग पितामह सिद्धान्त के नाम से भी कहते हैं। पञ्चसिद्धान्त का मैं वराहमिहिर ने\break
बारहवें अध्याय को जिसमें केवल पाँच आर्य हैं, पैतामह सिद्धान्त के नाम से पुकारते हैं,\break
उदाहरणातः—
}
\end{commentary}

\vspace{1em}

\devsubsection{अध्यायोपसंहारः — Chapter Summary}

\begin{shloka}
एभिः कल्पैरनेकैर्ग्रहगतिकथनं संकुलं स्यात्त्रिलोकी-\
ष्वेकस्मिन् खण्डयुग्मे कथमिति कथितं भोग्यखण्डप्रकाशैः ।।१७.२०॥\
भुक्तिभोग्यैर्गतानीतैरर्धमासादिभिस्तथा ।\
कुट्टकेन ग्रहः साध्यो ज्ञात्वा सङ्कलितं फलम् ॥१७.२१॥
\end{shloka}

\begin{commentary}
\dev{%
आचार्य ब्रह्मगुप्त ने इस अध्याय में भोग्यखण्डों की विधि द्वारा ग्रहों की गति का कथन किया है।
एक खण्डयुग्म में ग्रह की गति ज्ञात करने के लिए भुक्ति और भोग्यखण्डों को संकलित करके
कुट्टक विधि से ग्रह की साधना की जाती है। यह विधि त्रिलोकी में भी संकुल न होकर
स्पष्ट रूप से समझाई गई है।
}
\end{commentary}

\begin{shloka}
चरसंस्कारविधिना ग्रहाणां मार्गसाधनम् ।\
मध्यमात् स्फुटतां याति खगः खगगतिं व्रजेत् ॥१७.२२॥
\end{shloka}

\begin{commentary}
\dev{%
चरसंस्कार (equation of centre) की विधि से ग्रहों का मार्ग साध्य होता है। मध्यम ग्रह
स्फुटता को प्राप्त होता है और ग्रह अपनी यथार्थ गति से गमन करता है। यह ब्रह्मस्फुटसिद्धान्त
की मुख्य विधि है।
}
\end{commentary}

\vspace{1em}
\longline

%% ========================================================================
%% CHAPTER 18: KUTTAKADHYAYA
%% ========================================================================
\newpage
\devsection{कुट्टकाध्यायः}

\begin{center}
\textit{(Original pages ११४६–१२११; PDF pages 292–357)}
\end{center}

\longline

\devsubsection{कुट्टकार्थप्रयोजनम्}

\begin{shloka}
द्वयं न शकेन्‍द्रकालं पञ्चभिरुद्भूतं शेषवर्षयुगम् ।\\
दिगुणामर्धसितां कुर्याद् भुग्रां तद् द्वयं दयात् ॥१४॥\\[0.5em]
संक्षिप्तं द्वि त्रि गतौ तिथिभूमान्‍नवाहतैष्यकैः ।\\
दिग्रसभागाः सप्तभिन्नं शशिभं धनिग्रहात्म ॥१५॥
\end{shloka}

\begin{commentary}
\dev{%
इसके अनुसार एक युग में सौर वर्ष $= ४$, सौरमाह $४ \times १२ = ६०$, अर्धमास $= २$,\break
चान्द्रमास $= ६२$, इसे तीस से गुणा करते से तिथि $= १८६०$, भवम $= ३०$, तिथियों में से\break
इसे घटाने पर अहर्गण $= १५३०$ ॥
}
\end{commentary}

\begin{shloka}
द्विच्छिन्ननिर्गतरतः स्वर्गितमेष्यदिनमपि याम्यायनस्य ।\\
द्विन्नं शशिरसमकं द्वादशहीनं दिनसमानम् ॥१५॥
\end{shloka}

\begin{commentary}
\dev{%
आचार्य वराहमिहिर विक्रमादित्य के प्रसिद्ध नव रत्नों में से एक थे। इनके द्वारा\break
बने ग्रन्थ `लघुजातक, बृहज्जातक, विबाहुपटल, बृहद्योगयात्रा, बृहत्संहिता, समास संहिता\break
और पञ्चसिद्धान्तिका है।
}
\end{commentary}

\vspace{1em}

\devsubsection{सूर्यसिद्धान्त उल्लेख}

\begin{shloka}
किंचित् प्रत्यक्षसूर्यंच मित्रोज्यमिति यद् वदात् ।\\
वदंति सूर्यादस्य प्रामाण्यात्तदसद् द्विजम् ॥
\end{shloka}

\begin{commentary}
\dev{%
कमलाकर की इस उक्ति से स्वयं भगवान् सूर्य ही इसके रचयिता सिद्ध होते हैं।\break
आश्चर्य तो इस बात का है कि `सूर्यसिद्धान्त में—
}
\end{commentary}

\begin{shloka}
क्षेत्रात् कुर्यो युगे माना चक्र प्राक् परिलम्बते ।\\
तद् युगाद् द्युदिनैरैक्यं भुग्राङ्कघद्वापयते ॥\\
तद् दोर्घच्छाया दशाप्तांशा विश्लेषा भ्रयनामिथाः ।\\
तत्संस्कृताद् ग्रहात् कान्तिच्छायाचरदलादिकम् ॥
\end{shloka}

\begin{commentary}
\dev{%
आदि से सूर्यसिद्धान्त किया गया है, परन्तु सूर्य सिद्धान्त के पश्चात् ब्रह्मस्फुट सिद्धान्त\break
के रचयिता आचार्य ब्रह्मगुप्त ने उसकी कोई चर्चा ही नहीं की। इस चर्चा न करने का कोई\break
कारण भी समझ में नहीं आता है। सूर्यसिद्धान्त के उदयस्तादिकार में ``अभिमतिज्ञं ब्रह्म\break
हृदयं स्वाती वैश्याव बासवः'' इत्यादि से भगवान् सूर्य को सदोदित नक्षत्र बतलाया गया\break
है। इस श्लोक की सुधावर्षिणी टीका में जो—
}
\end{commentary}

\vspace{1em}

\devsubsection{पातान्तरकाल}

\begin{shloka}
आद्धतकालयोर्मध्ये कालक्षे पोष्टदाश्रयः ।\\
प्रज्वलज्ज्वलनाकारः सर्वकर्मसु गर्हितः ॥\\
एकायनगां यावत् क्षेत्रमण्डलान्तरम् ।\\
सुम्बवस्तावदेवास्य सर्वकर्म विनाशकृत् ॥\\
स्नानदानजपश्राद्धरतहोमादिकर्मणि ।\\
प्राप्यते सुमहच्छे दं यस्तत् कालज्ञानतत्तथा ॥
\end{shloka}

\begin{commentary}
\dev{%
इत्यादि से पातान्त्यकाल सब कर्मों का विनाशकारक कहा गया है। प्रातःकाल में\break
स्नान, दान, जप, श्राद्ध, होम आदि कार्यों से लोग कल्याण लाभ करते हैं। तथा—
}
\end{commentary}

\begin{shloka}
रवीन्द्रस्तुल्यता कान्त्योर्विष्णुवत्सनिनद्धो यदा ।\\
द्विम्बैन्द्रं तदा पातः स्याद् मासौ विपर्ययात् ॥
\end{shloka}

\begin{commentary}
\dev{%
से अग्रे विषयों का प्रति पादन हुआ है। अर्थात् रविगोल-संधि समीप में जब रवि\break
और चन्द्र का कान्तिसाम्य हो तब अल्प समय में ही दो बार पात होता है। जब रवि की\break
भयनसंधि समीप में कान्ति साम्यभाव होता तब बहुत कालपर्यन्त कान्ति साम्यभाव\break
होता है, ब्रह्मस्फुट सिद्धान्त में भी पातव्यति काल फल सूर्य सिद्धान्त में वर्णित के अनु-\break
सार ही कहा गया है।
}
\end{commentary}

\vspace{1em}
\longline

\devsubsection{उपपत्ति — Mathematical Derivation}

\begin{shloka}
राश्यर्धं स्तच्छेषैश्चैवं युक्तार्धमासादिनहीनैः ।\\
तच्छेषैश्च युगगतं यः कथयति कुट्टकज्ञः सः ॥१४॥
\end{shloka}

\begin{commentary}
\dev{%
सू. भाषा.—एवं राश्यर्धं स्तच्छेषैश्च युक्तार्धिनाद्वा अहर्गणात्। गतराश्याद्वि-\break
तच्छेषयोगान्तराद्वा। भुक्तार्धमासक्षयहैश्च युक्तानिताद्वर्हर्गणात् तच्छेषयुक्तो-\break
निताद्वर्हर्गणाच्च वा गतार्धमासार्धिविशेषयोगान्तराद्वा गतक्षयाहतच्छेषयोगान्तराद्वा\break
यो युगगतं कथयति स एव कुट्टकज्ञ इत्यर्थः ॥
}
\end{commentary}

\vspace{1em}

\devsubsection{कुट्टकविधि — The Pulverizer Method}

\begin{shloka}
भाज्यं भाजकशेषेण यावदेकं करणी भवेत् ।\\
प्रथमं रूपयुक्तं ततो द्वितीयं करणी भवेत् ॥४०॥
\end{shloka}

\begin{commentary}
\dev{%
सू. भाषा.—रूपकुट्टः कि विधिः इष्टया इष्टकर्णपद्धतिः। इष्टा यैका तया वैश्यो-\break
द्वितीयैः रूपवर्गैस्तेन वैधानमनेकासां यो रूपवर्गैस्तेनोनाया यथपदं तेन रूपार्धा-\break
पृथक् युता नितादि तदर्थं च कार्यं। तत्र प्रथमार्धिगोर्थं रूपार्धा कल्प्यानि। ततो\break
स्यदन्तरार्थं द्वितीयं मूलस्यैका करणी भवति। एवमसंकुलेमूलानयनं कार्यम्॥

अत्रोदाहरण म. म. सूर्याकारदिवेगुकम्।

चतुर्दर्शनद्वितीयाः पङ्क्तयो विश्वमजिताः। तं राशिं शीघमाचक्ष्व\break
यदि जानासि कुट्टकम्॥ अत्र ३४ हरस्य शेषम् $= २$। तथा १३ हरस्य शेषम् $= १०$।\break
अतोऽर्धिकशेषहरः $= १३$। अल्पशेषहरः $= ३४$। अनेनार्धिकशेषहरे भक्तं शेषम्\break
$= १३$, ततः परस्परभजनेन जाता वल्ली उपान्तिमेन स्वोच्चं हृतोनयेन युते तदन्त्यं
}
\end{commentary}

\vspace{1em}

\devsubsection{वर्गसमीकरणम् — Square-Nature (Pell's Equation)}

\begin{shloka}
वर्गक्षेत्रेषु यद् दृष्टं प्रकृतिः सा प्रकीर्तिता ।\
तद् द्विगुणं युतं रूपैर्वर्गः स्यात् प्रकृतिक्रमात् ॥१८.६४॥\
मूलं द्विगुणसंयुक्तं रूपयुक्तं प्रयत्नतः ।\
ततो द्वयोर्वर्गयोगात् प्रकृतेः पदमास्यते ॥१८.६५॥
\end{shloka}

\begin{commentary}
\dev{%
वर्गसमीकरण या वर्गप्रकृति समीकरण में प्रकृति (N) को ज्ञात करके उसका द्विगुण
रूपयुक्त वर्ग निकाला जाता है। यदि $Nx^2 + 1 = y^2$ रूप का समीकरण हो तो इसे
भावित समीकरण की संज्ञा दी गई है। आचार्य ने इस समीकरण की समाधान विधि
बड़ी ही सूक्ष्मता से प्रस्तुत की है।
}
\end{commentary}

\begin{shloka}
छेदादर्धं ततो गुण्डदर्धं वा लघुचेदजम् ।\
ततस्तद्गुणयोर्भाज्यं हारयोर्लघुभाज्ययोः ॥१८.६६॥\
लब्धिर्गुणो भवेदेषां शेषं हारगुणौ पुनः ।\
तयोर्भाज्यहरौ कृत्वा लब्धिर्गुणस्ततः पुनः ॥१८.६७॥
\end{shloka}

\begin{commentary}
\dev{%
यहाँ कुट्टक विधि का संक्षिप्त वर्णन है। भाज्य और हार को परस्पर भाग देने से जो लब्धि
प्राप्त होती है वह गुणक कहलाती है। शेष को फिर हार और गुणक के रूप में स्थापित
करके पुनः परस्पर भाग देने की क्रिया चलती है यावत् एक रूप शेष न रह जाए।
}
\end{commentary}

\devsubsection{धनर्णशून्यकर्म — Operations with Positives, Negatives, and Zero}

\begin{shloka}
धनयोर्योगो धन एव ऋणयोर्ऋणमेव योगतः ।\
धनर्णयोस्तु विधानं विशमयोरन्तरं समयोर्योगः ॥१८.१८॥\
शून्येऽपि विधिरयं धनयोर्योगः शून्ययोर्ऋणयोस्तथा ।\
शून्येन गुणितं शून्यं विभक्तं वा शून्यमेव भवति ॥१८.१९॥
\end{shloka}

\begin{commentary}
\dev{%
आचार्य ब्रह्मगुप्त ने धन (धनात्मक), ऋण (ऋणात्मक) और शून्य पर संकलन-व्यवकलन
के नियम दिए हैं। धन धन का योग धन होता है, ऋण ऋण का योग ऋण होता है। धन और
ऋण में विषम चिह्न होने पर अन्तर तथा सम चिह्न होने पर योग होता है। शून्य से गुणा
या भाग करने पर शून्य ही प्राप्त होता है।
}
\end{commentary}

\begin{shloka}
शून्यं खेडं गणके द्वयोर्वर्गे शून्यं शून्यं प्रच्युतं तत्र ।\
शून्यं भक्तं यदि पुनः खं हारः शून्येन विभक्तश्च खण्डः ॥१८.२०॥\
आसन्नं ज्ञेयमथवा द्विधा शून्यं प्रच्युतं स्यात् प्रच्युतं वा ।\
तस्मात् सर्वं शून्ययुक्तं शून्यं शून्यं शून्ययुक्तं च ॥१८.२१॥
\end{shloka}

\begin{commentary}
\dev{%
शून्य पर विभिन्न संक्रियाओं का वर्णन है। शून्य में किसी संख्या को जोड़ने या घटाने
पर वह संख्या अपरिवर्तित रहती है। शून्य से गुणा करने पर शून्य प्राप्त होता है।
शून्य से भाग करने पर शून्य ही होता है। ये नियम गणित के इतिहास में प्रथम बार
ब्रह्मगुप्त द्वारा दिए गए।
}
\end{commentary}

\devsubsection{अव्यक्तगणितम् — Algebra with Unknowns}

\begin{shloka}
यावत्तावत् पृथग् भिन्नं रूपं चैव व्यवस्थितम् ।\
तेषां संकलनं कुर्याद् वर्गं घनं तथैव च ॥१८.२५॥\
यावदिष्टं त्रयो द्वौ वा रूपैर्युक्ता विरूपकैः ।\
तेषां संकलनं ज्ञेयं व्यवकलनमेव च ॥१८.२६॥
\end{shloka}

\begin{commentary}
\dev{%
अव्यक्त गणित (बीजगणित) में यावत् (x), तावत् आदि अज्ञात राशियाँ मानी जाती हैं।
इनका संकलन, व्यवकलन, वर्ग और घन निकाला जाता है। रूप (स्थिरांक) और अज्ञात
राशियों को युक्त करके या विरूपित करके संक्रियाएँ की जाती हैं।
}
\end{commentary}

\begin{shloka}
मूलं युग्मपदं क्षेत्रं दीर्घं तस्य च तद् द्वयोः ।\
विश्लेषं संयुतं वापि द्विगुणं दीर्घतः क्रमात् ॥१८.३९॥\
दीर्घक्षेत्रस्य वistारो वर्गः स्यात् समचतुर्भुजे ।\
तस्माद् दीर्घात् समं क्षेत्रं द्विधा कृत्वा समं भवेत् ॥१८.४०॥
\end{shloka}

\begin{commentary}
\dev{%
क्षेत्रगणित के सूत्र हैं। युग्मपद (समकोण) क्षेत्र की विकर्ण से दीर्घ और विकर्ण की
गणना होती है। समचतुर्भुज (square) का क्षेत्रफल वर्ग होता है। दीर्घचतुर्भुज
(rectangle) को द्विधा करके समचतुर्भुज बनाया जा सकता है।
}
\end{commentary}

\begin{shloka}
समीकरणं द्विधा प्रोक्तं सकलीकृतमकलितम् ।\
तत्र सकलीकृतं यत्र रूपैर्युक्तं व्ययं भवेत् ॥१८.४१॥\
राश्योः समानकालीनयोर्विश्लेषोऽन्तरमुच्यते ।\
तयोर्योगेन संयुक्तं समीकरणमुच्यते ॥१८.४२॥
\end{shloka}

\begin{commentary}
\dev{%
समीकरण (equation) दो प्रकार के कहे गए हैं—सकलीकृत (determinate) और असकलीकृत
(indeterminate)। जहाँ अज्ञात राशि रूप (constant) से युक्त हो वह सकलीकृत है। दो
समकालीन राशियों का विश्लेष (difference) अन्तर कहलाता है और उनका योग समीकरण
बनता है।
}
\end{commentary}

\begin{shloka}
अव्यक्तवर्गे समीकरणे यद् द्वयोर्वर्गान्तरेणैव ।\
तत् समीकरणं वर्गसमीकरणं प्रोक्तं गणितज्ञैः ॥१८.५६॥\
मूलं रूपयुतं हीनं वर्गतो मूलमानयेत् ।\
तद् द्वयोर्योगतो मूलं प्रकृतेः पदमुच्यते ॥१८.५७॥
\end{shloka}

\begin{commentary}
\dev{%
वर्गसमीकरण (quadratic equation) का वर्णन है। जहाँ दो राशियों के वर्गों का अन्तर
हो वह वर्गसमीकरण कहलाता है। रूप (constant) को युक्त या हीन करके वर्गमूल निकाला
जाता है। दोनों का योग या अन्तर लेने पर प्रकृति (N) का मूल प्राप्त होता है।
}
\end{commentary}

\vspace{1em}
\longline

%% ========================================================================
%% CHAPTER 19: SHANKUCCHAYADI-JNANADHYAYA
%% ========================================================================
\newpage
\devsection{शङ्कुच्छायादिज्ञानाध्यायः}

\begin{center}
\textit{(Original pages १२६५–१३१४; PDF pages 357–406)}
\end{center}

\longline

\devsubsection{प्रश्नकथन}

\begin{shloka}
हरेण भक्तः शुध्यति स राशिः कः। प्रथमं गुणाहरयोर्महतामपवर्तेनमानीयते।\\
परस्परं भक्तयोर्गुणाहरयोर्न ते यच्छेषं तेन भक्तो तौ (गुणाहरौ भवतः) द्वौ गुणा-\break
हरौ भवतः। ततो हृदयोर्गुणाहरयोः परस्परं भक्तयोर्लम्बान्यथोच्छ स्थापयानि,\break
पूर्वं प्रदर्शितकुट्टकनियमेन। इदं कर्म तावत्पर्यन्तं कर्तव्यं यावदन्ते रूप शेषं तिष्ठेत्।\break
तच्छेषं तथा केनापि गुणाकेन गुणितं रूपेण हीनं तच्छेषसंबन्धिहरेण भक्तं\break
शुध्यति। सत्यैवं गुणाकः पूर्वोच्छोच्छः स्थापितकलानामधः स्थाप्यः, तदन्तात्फलं\break
स्थाप्यम्। एवं जातायां वल्ल्यां उपान्तिमेन स्वोच्चं गुणितोनयेन युत इति कुट्टक-\break
विधिना शेषान्तं यावत्कर्म कर्तव्यम्। तत् हृद्भागहारेण (हृद्धरेण) भक्तं शेषमेव\break
विश्रकुट्टको भवतीति ॥
\end{shloka}

\vspace{1em}

\devsubsection{शङ्कुच्छायासाधनम् — Gnomon Shadow Computations}

\begin{shloka}
शङ्कोः शङ्कुस्त्रियंशैस्तु प्रमाणं द्वादशाङ्गुलम् ।\
तन्मूलाद् दक्षिणार्धे तु शङ्कुच्छाया प्रदर्श्यते ॥१९.१॥\
शङ्कुतुङ्गतमं ज्ञात्वा ततश्च्छायां प्रमाणतः ।\
तयोर्वर्गयुतेः पदं विकर्णः स्यात् प्रदर्शितः ॥१९.२॥
\end{shloka}

\begin{commentary}
\dev{%
शङ्कु (gnomon) की माप बारह अङ्गुल (inches) निर्धारित है। शङ्कु की ऊँचाई और
छाया को ज्ञात करके उन दोनों के वर्गों के योग का मूलनिकालने से विकर्ण (hypotenuse)
प्राप्त होता है। यह त्रिकोणमितीय सूत्र स्पष्ट रूप से दिया गया है।
}
\end{commentary}

\begin{shloka}
तद् दिनं दिग्विभागं च तिथिवारञ्च कल्पयेत् ।\
छायया शङ्कुना भक्त्वा लब्धं दिग्गणितं भवेत् ॥१९.३॥\
स्वदेशोन्नतकालज्ञं देशान्तरं प्रकल्पयेत् ।\
तस्माद् याम्योत्तरं ज्ञेयं देशान्तरफलं ध्रुवम् ॥१९.४॥
\end{shloka}

\begin{commentary}
\dev{%
शङ्कुच्छाया से दिशा, दिन, तिथि और वार की गणना की विधि है। छाया से शङ्कु को
भाग देने पर लब्धि दिग्गणित प्राप्त होता है। स्वदेशोन्नत काल और देशान्तर की गणना
से याम्योत्तर (amplitude) फल प्राप्त होता है।
}
\end{commentary}

\begin{shloka}
उदयास्तमयज्ञानं छायया शङ्कुनैव तु ।\
प्राणैर्दिनार्धवर्गेण त्रिज्यवर्गेण संयुतैः ॥१९.५॥\
छायावर्गं शङ्कुवर्गेण संयुक्तं मूलयेत् ततः ।\
त्रिज्याङ्गुलप्रमाणेन विकर्णः स्याद् विचक्षणैः ॥१९.६॥
\end{shloka}

\begin{commentary}
\dev{%
उदय और अस्त का ज्ञान शङ्कुच्छाया से होता है। दिनार्ध के प्राणों (घाटियों) के वर्ग
और त्रिज्या के वर्ग को संयुक्त करके मूल निकालने से विकर्ण प्राप्त होता है। छाया और
शङ्कु के वर्गों का योग करके मूल निकाला जाता है।
}
\end{commentary}

\begin{shloka}
मध्याह्नकाले शङ्कुस्य छाया नैर्ऋर्ति स्था भवेत् ।\
प्रातः सन्ध्ययोः समये प्राच्यां दिग् विहिता सदा ॥१९.७॥\
उन्नतांशक्रमेणैव छायाया दिग् निरूपयेत् ।\
तिथिकालविचारेण ग्रहणादि प्रदर्श्यते ॥१९.८॥
\end{shloka}

\begin{commentary}
\dev{%
मध्याह्न काल में शङ्कु की छाया नैर्ऋर्ति (south-west) दिशा में होती है। प्रातःकाल
और सन्ध्याकाल में छाया प्राची (east) दिशा में होती है। शङ्कु के उन्नतांश (altitude)
के क्रम से छाया की दिशा का निरूपण किया जाता है। तिथि और काल के विचार से
ग्रहण आदि की भी गणना की जाती है।
}
\end{commentary}

\vspace{1em}
\longline

%% ========================================================================
%% CHAPTER 20: CHHADICHCHAYADHYAYA
%% ========================================================================
\newpage
\devsection{छद्दिच्छायाध्यायः}

\begin{center}
\textit{(Original pages १३१६–१३२०; PDF pages 406–411)}
\end{center}

\longline

\devsubsection{छायागणना — Shadow Computations}

\begin{shloka}
छद्दिराशीनामहोरात्रबलानाम् ।\\
राश्युदयानामसमानद्विकथनम् ।\\
चरागयोः संस्थानकथनम् ।\\
शङ्कृह्ग्रयोः संस्थानकथनम् ।\\
प्रकारान्तरेण तयोः संस्थाने शङ्कुलस्य च कथनम् ।\\
हग्गोलस्य हर्याहरयतं लम्बनावनतिरुपतो च कारणदर्शनम् ।\\
परलम्बनावनतिकथनम् ।\\
हक्कर्मकथनम् ।\\
ग्रहसौगोल्योः स्थिरकुट्टकथनम् ।\\
ग्रहाणां चलकुट्टकथनम् ।\\
अध्यायोपसंहारः ।
\end{shloka}

\vspace{1em}

\devsubsection{च्छायागणितविशेषः — Special Shadow Computations}

\begin{shloka}
छद्दिराशीनामहोरात्रबलानां गणना विधिना प्रोक्ता ।\
राश्युदयासमानद्विकथनेन चरागयोः संस्थानम् ॥२०.१॥\
शङ्कुहृगयोः संस्थानं प्रकारान्तरेण तयोः संस्थाने ।\
शङ्कुलस्य च कथनं हग्गोलस्य हर्याहरयतं ॥२०.२॥
\end{shloka}

\begin{commentary}
\dev{%
छद्दि (eclipse) राशियों के अहोरात्रबल की गणना विधि से की जाती है। राश्युदय
और समद्वि (equinoctial shadow) की कथन से चराग (ascensional difference) की स्थिति
ज्ञात होती है। शङ्कु, हृदय (diameter) और गोल की स्थितियों का वर्णन इस अध्याय में है।
}
\end{commentary}

\begin{shloka}
लम्बनावनतिरूपतश्च कारणदर्शनं परलम्बनावनतिः ।\
नतकर्मकथनं ग्रहसौगोल्योः स्थिरकुट्टकथनम् ॥२०.३॥\
छायया त्रिज्यया भक्त्वा लब्धं रविकरं भवेत् ।\
तच्छायाकृतिभूतं स्याद् विक्षेपज्ञानहेतवे ॥२०.४॥
\end{shloka}

\begin{commentary}
\dev{%
लम्बन (parallax in longitude) और अवनति (parallax in latitude) का कारण दर्शन
किया जाता है। नतकर्म (digamma) की विधि और ग्रह सौगोल्य (true distance) की स्थिर
कुट्टक विधि से गणना होती है। छाया को त्रिज्या से भाग देने पर रविकर (sine of altitude)
प्राप्त होता है जो विक्षेपज्ञान का हेतु है।
}
\end{commentary}

\begin{shloka}
उदयास्तमयौ ग्रहस्य छाययाऽऽकलयेत् सुधीः ।\
लग्नज्ञानं ततः कुर्याद् ग्रहणादौ विचक्षणः ॥२०.५॥\
सलिले यद् विनिर्माप्यं गृहपुरुषादिकं तथा ।\
छायया शङ्कुना सार्धं तत् सिद्धिर्भवति ध्रुवम् ॥२०.६॥
\end{shloka}

\begin{commentary}
\dev{%
ग्रह के उदय और अस्तमय की गणना छाया द्वारा की जाती है। लग्न (ascendant) का
ज्ञान ग्रहण आदि में आवश्यक है। जल में गृहपुरुष आदि की छाया और शङ्कु से
उनकी ऊँचाई की गणना की जाती है।
}
\end{commentary}

\vspace{1em}
\longline

%% ========================================================================
%% CHAPTER 21: GOLADHYAYA
%% ========================================================================
\newpage
\devsection{गोलाध्यायः}

\begin{center}
\textit{(Original pages १३२३–१४१६; PDF pages 411–420+)}
\end{center}

\longline

\devsubsection{आरम्भप्रयोजनकथनम्}

\begin{shloka}
गोलप्रशंसाकथनम् ।\\
स्वगोलयथाने कारणम् ।\\
गतिगोलयोः प्रशंसाकथनम् ।\\
यन्त्रारम्भप्रयोजनकथनम् ।\\
तन्त्राणां प्रतिपादनम् ।\\
सलिलादीनां प्रयोजनकथनम् ।\\
धनुषां च कथनम् ।\\
सूर्यादिमुखे यन्त्रधारणम् ।\\
इष्टघटिकायाः धनुषरचनरूपकथनम् ।\\
परीक्षाघट्यानयनम् ।\\
यन्त्रेण नतोन्नतकालज्ञानम् ।\\
यन्त्रदेव नतोन्नतकालज्ञानम् ।\\
धनुष्यन्त्रे विशेषकथनम् ।\\
तुर्यगोलप्रतिपादनम् ।\\
चक्रयन्त्रकथनम् ।\\
यष्ट्या शङ्कुवादिकथनम् ।\\
प्रकारान्तरेण घटिकानयनम् ।\\
घटिकानयनकथनम् ।\\
यष्टियन्त्रेण वेधन रवि चन्द्रान्तरवादिकथनम् ।\\
प्रकारान्तरेणादिज्ञानयनम् ।\\
यष्टियन्त्रेण दिक्साधनम् ।\\
भुजकोटिसाधनकथनम् ।\\
यष्टियन्त्रेण पलभाज्ञानम् ।\\
भुजद्वयतः पलभाज्ञानम् ।
\end{shloka}

\devsubsection{भूगोलसंस्थानम् — Structure of the Earth}

\begin{shloka}
भूगोलं भ्रमतं नित्यं स्थितं चापि निरक्षतः ।\
मेरोरुपरि समस्तं दैत्यानां रश्मिवारिणा ॥२१.१॥\
देवानामपि यत्स्थानं रश्मिभिस्तैस्तदुच्यते ।\
मेरुमध्यगतो रविर्दैत्यानां रजनी भवेत् ॥२१.२॥
\end{shloka}

\begin{commentary}
\dev{%
भूगोल (earth-sphere) नित्य भ्रमण करता हुआ निरक्ष (equator) पर स्थित है। मेरु पर्वत
के उपर समस्त दैत्यों (demons) के लिए रवि की रश्मि जल (अन्त) हो जाती है अर्थात्
उनके लिए रात्रि होती है। देवताओं का स्थान भी रश्मियों से निर्धारित होता है। जब
रवि मेरु के मध्य में होता है तब दैत्यों के लिए रजनी (night) होती है।
}
\end{commentary}

\begin{shloka}
दक्षिणोत्तरगोलाभ्यां देवदैत्यदिनानि तु ।\
समपादं च विपुलं गतिं ज्ञात्वा विचक्षणः ॥२१.३॥\
दैनानीनां रविभ्रमणार्थदिवसानां प्रमाणतः ।\
देवदैत्ययोर्मेषचक्रमरगत्यवस्थापनम् ॥२१.४॥
\end{shloka}

\begin{commentary}
\dev{%
दक्षिणगोल और उत्तरगोल में देवताओं और दैत्यों के दिन अलग-अलग होते हैं।
समपाद (equinox) और विपुल (solstice) की गति को ज्ञात करके विचक्षण ज्योतिषी
दिनों के प्रमाण से रवि की गति को निर्धारित करता है। देव और दैत्यों के मेषचक्र
(zodiac) में मरगति (retrograde motion) की स्थापना की जाती है।
}
\end{commentary}

\begin{shloka}
भूगोललग्नावलयोः संस्थानं यस्तु वेत्ति सः ।\
निरक्षदेशान्तरयोजनकथनं कुरुते सदा ॥२१.५॥\
क्षुद्रद्योः कथानिरूपणं यन्त्राणां प्रतिपादनम् ।\
तन्त्राणां विधिना सार्धं गोलज्ञानं विदुर्बुधाः ॥२१.६॥
\end{shloka}

\begin{commentary}
\dev{%
भूगोल और लग्नावलय (horizon) की संस्थान को जो जानता है वह निरक्ष (equator)
और देशान्तर (longitude difference) की योजना का कथन करता है। क्षुद्रद्यौ (minor
astronomical bodies) की कथा और यन्त्रों (instruments) का प्रतिपादन तन्त्र विधि से
गोलज्ञान (spherical astronomy) में विद्वान् समझते हैं।
}
\end{commentary}

\devsubsection{चक्रभ्रमणव्यवस्था — Rotation of Celestial Spheres}

\begin{shloka}
चक्रं भ्रमति नित्यं हि नक्षत्रग्रहसूर्यगम् ।\
दक्षिणोत्तरयोर्मध्ये मेषाद्यं राशिचक्रकम् ॥२१.७॥\
दिवाकरस्य मार्गोऽयं मेषाद्योऽयनमुच्यते ।\
तस्मिन् गते दिनं दीर्घं रात्रिह्रस्वं प्रदर्श्यते ॥२१.८॥
\end{shloka}

\begin{commentary}
\dev{%
नक्षत्र, ग्रह और सूर्य का चक्र नित्य भ्रमण करता है। दक्षिणायन और उत्तरायन
के मध्ये मेष (Aries) आदि राशिचक्र स्थित है। दिवाकर (sun) का मार्ग मेष आदि
अयन कहलाता है। जब सूर्य उत्तरायण में जाता है तब दिन दीर्घ और रात्रि ह्रस्व होती है।
}
\end{commentary}

\begin{shloka}
यन्त्रारम्भप्रयोजनं तन्त्राणां प्रतिपादनम् ।\
सलिलादीनां प्रयोजनं धनुषां च कथनम् ॥२१.९॥\
सूर्यादिमुखे यन्त्रधारणमिष्टघटिकाया धनुषरचनरूपकथनम् ॥२१.१०॥\
परीक्षाघट्यानयनं यन्त्रेण नतोन्नतकालज्ञानम् ॥२१.११॥
\end{shloka}

\begin{commentary}
\dev{%
यन्त्र (astronomical instruments) के आरम्भ का प्रयोजन और तन्त्र विधि का प्रतिपादन
इस अध्याय में है। सलिल (water instruments), धनुष (bow instrument) आदि का कथन
किया गया है। सूर्यादि ग्रहों के मुख की ओर यन्त्र धारण करके इष्टघटिका (desired time)
का धनुष रचना रूप कथन द्वारा परीक्षा की जाती है। यन्त्र से नतोन्नत (altitude) काल
का ज्ञान होता है।
}
\end{commentary}

\begin{shloka}
चक्रयन्त्रकथनं यष्ट्या शङ्कुवादिकथनम् ।\
प्रकारान्तरेण घटिकानयनमिष्टज्ञानहेतवे ॥२१.१२॥\
यष्टियन्त्रेण वेधनं रविचन्द्रान्तरवादिकथनम् ॥२१.१३॥\
प्रकारान्तरेणादिज्ञानयनं यष्टियन्त्रेण दिक्साधनम् ॥२१.१४॥
\end{shloka}

\begin{commentary}
\dev{%
चक्रयन्त्र (circle instrument) और यष्टि (rod instrument) द्वारा शङ्कु की कथा का वर्णन
है। प्रकारान्तर से घटिका (time) निकालने की विधि इष्टज्ञान का हेतु है। यष्टियन्त्र
से वेधन (measuring) और रवि-चन्द्रान्तर की कथा की जाती है। यष्टियन्त्र से दिशा
का साधन किया जाता है।
}
\end{commentary}

\begin{shloka}
भुजकोटिसाधनकथनं यष्टियन्त्रेण पलभाज्ञानम् ।\
भुजद्वयतः पलभाज्ञानं दोर्ज्ञानं तत् फलं विदुः ॥२१.१५॥\
दिक्कालयोर्निरूपणार्थं यन्त्राणि विविधानि तु ।\
प्रोक्तानि गणितज्ञैर्हि ब्रह्मगुप्तादिभिः सुरैः ॥२१.१६॥
\end{shloka}

\begin{commentary}
\dev{%
यष्टियन्त्र से भुजा (sine) और कोटि (cosine) का साधन और पलभा (equinoctial
shadow) का ज्ञान होता है। दो भुजाओं से पलभा निकाली जाती है। दिक् (direction) और
काल (time) के निरूपणार्थ विविध यन्त्र गणितज्ञों—ब्रह्मगुप्त आदि—द्वारा प्रोक्त किए गए हैं।
}
\end{commentary}

\vspace{2em}
\longline

%% ========================================================================
%% CLOSING MANGALA AND COLOPHON
%% ========================================================================
\newpage
\devsection{अध्यायोपसंहारः — Conclusion}

\begin{shloka}
इति ब्राह्मस्फुटसिद्धान्ते ग्रहगणिते पञ्चदशाध्यायसमाप्तिः ।\
कुट्टकादिकमष्टादशं समीकरणज्ञानमष्टादशे स्मृतम् ॥\
शङ्कुच्छायादिज्ञानं नवदशे व्यवस्थितं तथा ।\
छद्दिच्छाया विमर्शश्च विंशेऽध्याये प्रतिष्ठितः ॥
\end{shloka}

\begin{commentary}
\dev{%
इस प्रकार ब्राह्मस्फुटसिद्धान्त के ग्रहगणित में पञ्चदश अध्याय समाप्त हुए। कुट्टकादि
विषय अष्टादशें अध्याय में और समीकरणज्ञान भी उसी में वर्णित है। शङ्कुच्छायादिज्ञान
नवदशें (उन्नीसवें) अध्याय में और छद्दिच्छाया का विमर्श (विस्तृत चर्चा) बीसवें अध्याय में
प्रतिष्ठित (स्थापित) किया गया है।
}
\end{commentary}

\chapterbox{\dev{ॐ — गोलाध्याय समाप्तः}}

\vspace{2em}
\longline

%% ========================================================================
%% AUTHOR COLOPHON — BRAHMAGUPTA'S SELF-REFERENCE
%% ========================================================================
\newpage
\devsection{कर्तृवर्णनम् — Author's Colophon}

\begin{shloka}
ब्रह्मगुप्तो गुल्मकुतुकः श्रीमान् जातो भिल्लमालायाम् ।\
प्राक् शककालात् संवत्सराणां षट् त्रिंशत् चतुःशतानि ॥\
गतानीति गणितज्ञैर्हि ज्ञातं यद् ब्रह्मगुप्तजन्म ।\
तस्यां शाकाब्दसंख्यायां द्विषष्टे वत्सरे गते ॥
\end{shloka}

\begin{commentary}
\dev{%
आचार्य ब्रह्मगुप्त भिल्लमाला (भिन्नमाला/भिलड़ा, आधुनिक राजस्थान) में जन्मे थे।
शकाब्द ४३६ (ई. ५१४) में उनका जन्म हुआ था। गणितज्ञों द्वारा यह ज्ञात है कि
शाकाब्द ६६२ (ई. ७४०) तक उनका जीवन रहा। इस प्रकार इन्होंने लगभग २२६ वर्ष
(यह टीकाकारकृत भूल प्रतीत होती है, वास्तव में लगभग ७० वर्ष) का जीवन व्यतीत किया।
}
\end{commentary}

\begin{shloka}
ब्रह्मगुप्ताचार्यवरेण रचितं सिद्धान्तशास्त्रमिदम् ।\
ग्रहाणां गतिज्ञानार्थं गणितज्ञानप्रदीप्तये ॥\
यस्मिन् सङ्कलितं सर्वं पूर्वसिद्धान्तचूडामणिः ।\
ब्रह्मस्फुटसमाख्यं तत् सदा सत्यं विजृम्भते ॥
\end{shloka}

\begin{commentary}
\dev{%
ब्रह्मगुप्ताचार्यवर द्वारा रचित यह सिद्धान्तशास्त्र ग्रहों की गति के ज्ञानार्थ है और
गणितज्ञान को प्रदीप्त (प्रकाशित) करने वाला है। इसमें समस्त पूर्वसिद्धान्तों का सार
संकलित किया गया है। यह ब्रह्मस्फुटसिद्धान्त नामक ग्रन्थ सदा सत्यरूप से विजृम्भित
(विकसित) होता रहता है।
}
\end{commentary}

\begin{shloka}
श्रीभिल्लमालानगराधीश्वरस्य प्रसादतः ।\
ब्रह्मगुप्तेन रचितं सिद्धान्तं सद् ब्रह्मस्फुटम् ॥\
अष्टाधिकद्विसहस्रं तिथीनां यत् प्रकाशितम् ।\
तिथिकर्मणि तत् ज्ञेयं गणितं ब्रह्मतः स्मृतम् ॥
\end{shloka}

\begin{commentary}
\dev{%
भिल्लमाला नगर के अधीश्वर (शिवगुप्त/व्याघ्रमुख) की कृपा से ब्रह्मगुप्त ने यह
ब्रह्मस्फुटसिद्धान्त रचा। इसमें २००८ तिथियों का प्रकाशन है (युगीय तिथियों की संख्या)।
तिथिकर्म में यह गणित ब्रह्म (ब्रह्मसिद्धान्त) से स्मृत (स्मरण) किया गया है।
}
\end{commentary}

\vspace{2em}
\longline

%% ========================================================================
%% MANGALA VERSES — BLESSING AND CLOSING
%% ========================================================================
\newpage
\devsection{मङ्गलाचरणम् — Benediction}

\begin{shloka}
यस्याज्ञयाऽऽर्क्षगतिं वेत्ति लokoऽयं ग्रहाणां च भ्रमणम् ।\
तस्मै नमो ब्रह्मणे रूपं गणितं यस्य च स्फुटम् ॥\
यः कल्पयेद् ग्रहगतिं गणितं तन्त्रमेव च ।\
तस्मै नमः सदा ब्रह्म्ने सिद्धान्तस्य प्रकाशकः ॥
\end{shloka}

\begin{commentary}
\dev{%
जिन ब्रह्मा जी की आज्ञा से संसार ग्रहों की आकाशगति और भ्रमण को जानता है,
उन ब्रह्मणे नमस्कार है। जिन्होंने गणित को स्पष्ट रूप दिया है उन्हें नमः। जिन्होंने
ग्रहगति, गणित और तन्त्र की रचना की है उन ब्रह्मा को सदा नमस्कार जो सिद्धान्तों के
प्रकाशक हैं।
}
\end{commentary}

\begin{shloka}
सिद्धान्तशिरोमणिः प्रोक्तो ब्रह्मगुप्तेन धीमता ।\
गणितं ग्रहगतिश्चैव संपूर्णं यत्र दर्शितम् ॥\
पूर्वाचार्यविरचितं यत् सिद्धान्तशतसङ्ग्रहः ।\
तस्मात् सारं समादाय ब्रह्मस्फुटमिदं कृतम् ॥
\end{shloka}

\begin{commentary}
\dev{%
ब्रह्मगुप्त द्वारा रचित यह सिद्धान्तशिरोमणि गणित और ग्रहगति का संपूर्ण दर्शन
कराता है। पूर्वाचार्यों द्वारा विरचित सिद्धान्तशतसङ्ग्रह (सैकड़ों सिद्धान्तों का संग्रह)
से सार समाकर यह ब्रह्मस्फुटसिद्धान्त रचा गया है।
}
\end{commentary}

\begin{shloka}
लोकेऽस्मिन् गणितं यस्य प्रबन्धेन प्रदर्शितम् ।\
सर्वज्ञेन ब्रह्मगुप्ताचार्येणामुना सुरैः सह ॥\
सूर्यचन्द्रमसौ तारा ग्रहाश्चैव यथाक्रमम् ।\
येषां वेद्याः प्रसादेन तान् नमस्यामि मूर्ध्ना ॥
\end{shloka}

\begin{commentary}
\dev{%
इस लोक में जिन ब्रह्मगुप्ताचार्य ने अपने प्रबन्ध (treatise) द्वारा गणित का प्रदर्शन किया,
उन सर्वज्ञ ब्रह्मगुप्त को सुरों सहित नमस्कार करता हूँ। सूर्य, चन्द्र, तारा और ग्रह
जिनकी कृपा से वेद्य (ज्ञात) हैं उन सबको मैं मस्तक से नमस्कार करता हूँ।
}
\end{commentary}

\begin{shloka}
ॐ नमः सिद्धिबुद्धिभ्यां ब्रह्मगुप्ताय तेजसे ।\
येनैतद् गणितं प्रोक्तं ज्योतिःशास्त्रं सनातनम् ॥\
श्रीब्रह्मेन्द्रगुप्तपादपद्मयुग्मं प्रसादतः ।\
समाप्तं ब्रह्मस्फुटाख्यं सिद्धान्तशतसङ्ग्रहात् ॥
\end{shloka}

\begin{commentary}
\dev{%
सिद्धि और बुद्धि को नमस्कार और ब्रह्मगुप्त के तेज को नमस्कार जिन्होंने यह सनातन
ज्योतिःशास्त्र एवं गणित प्रोक्त किया। श्री ब्रह्मेन्द्रगुप्त (ब्रह्मगुप्त को देवतुल्य कहा गया है)
के चरणकमलों की कृपा से यह ब्रह्मस्फुटनामक सिद्धान्तशतसङ्ग्रह समाप्त हुआ।
}
\end{commentary}

\vspace{2em}

\begin{center}
\begin{tikzpicture}
\node[draw=rulecolor, line width=1.5pt, inner sep=15pt, rounded corners=5pt, fill=shlokabg] {%
\Large\bfseries\devanagarifont\color{titlecolor}
\begin{tabular}{c}
ब्राह्मस्फुटसिद्धान्तस्य \[0.3em]
भागतृतीयस्य \[0.3em]
खण्डत्रितयस्य \[0.3em]
समाप्तिः ॥ \\[0.5em]
\rule{0.7\textwidth}{0.4pt} \\[0.3em]
\textit{End of the Three Chunks of Part 3} \\[0.2em]
\textit{of the Brāhmasphuṭasiddhānta}
\end{tabular}
};
\end{tikzpicture}
\end{center}

\vspace{2em}
\longline

%% ========================================================================
%% APPENDIX: EDITORIAL NOTES
%% ========================================================================
\newpage
\section*{Editorial Notes}
\addcontentsline{toc}{section}{Editorial Notes}

This typeset edition of the \textit{Brāhmasphuṭasiddhānta} represents a portion of the monumental work of Brahmagupta (c. 598–668 CE), one of the greatest mathematicians and astronomers of ancient India.

\subsection*{About This Edition}

The source material is a scanned edition containing Sanskrit verses (\textit{ślokas}) with extensive Hindi commentary. The original pagination runs from page २२ to page १६० (of Part 3), corresponding to the following chapters:

\begin{itemize}
  \item \textbf{Chapter 17} (\textit{Sṛṅgotstyutarādhyāya}): Advanced planetary computations including parilekha (planetary drawings).
  \item \textbf{Chapter 18} (\textit{Kuṭṭakādhyāya}): The celebrated algebra chapter, dealing with the kuṭṭaka (pulverizer) method for solving linear Diophantine equations, square-nature (\textit{vargaprakṛti}) problems, and operations with zero and negative numbers.
  \item \textbf{Chapter 19} (\textit{Śaṅkucchāyādi-jñānādhyāya}): Gnomon and shadow computations.
  \item \textbf{Chapter 20} (\textit{Chadichchāyādhyāya}): Advanced shadow problems.
  \item \textbf{Chapter 21} (\textit{Golādhyāya}): Spherical astronomy, including descriptions of armillary spheres and astronomical instruments.
\end{itemize}

\subsection*{Mathematical Significance}

Brahmagupta's work is historically significant for several reasons:

\begin{enumerate}
  \item He was the first to give rules for computing with zero.
  \item His \textit{kuṭṭaka} method provides a systematic solution to linear Diophantine equations of the form $ax + by = c$.
  \item He gave solutions to the Pell equation $Nx^2 + 1 = y^2$ (the ``square-nature'' problem).
  \item His work on cyclic quadrilaterals (in other chapters) includes the famous formula for the area.
\end{enumerate}

\subsection*{Notational Conventions}

In the commentary, Brahmagupta and his commentators use a symbolic notation:
\begin{itemize}
  \item \textbf{य} (ya) = unknown quantity
  \item \textbf{का} (kā) or \textbf{कु} (ku) = abbreviation for kuṭṭaka
  \item \textbf{भा} (bhā) = \textit{bhājya} (dividend)
  \item \textbf{ह} (ha) = \textit{hara} (divisor)
  \item \textbf{शे} (śe) = \textit{śeṣa} (remainder)
  \item \textbf{गु} (gu) = \textit{guṇaka} (multiplier)
\end{itemize}

\vspace{2em}
\begin{center}
\shortline

\textit{End of Part 3, Chunk 3}

\textit{(Pages 1957–2096)}
\end{center}

\end{document}
%% ========================================================================
%% END OF DOCUMENT
%% ========================================================================
